\chapter{Echocardiogram}

\section*{Overview}
An echocardiogram is an ultrasound examination that creates moving images of the heart and uses Doppler measurements to assess blood flow. The common transthoracic study scans through the chest; transesophageal, stress, and contrast-enhanced studies answer different clinical questions.

\section*{Why It Is Done}
It evaluates murmurs, heart failure, valve disease, cardiomyopathy, congenital heart disease, pericardial fluid, and suspected structural or functional abnormalities.

\section*{How It Is Performed}
For a transthoracic study, a probe is moved over gelled areas of the chest. A transesophageal study passes a specialized probe into the esophagus after sedation to obtain closer images when needed.

\section*{Preparation}
Standard transthoracic echocardiography usually requires little preparation. A transesophageal study has separate fasting and sedation instructions. Tell the team if you have difficulty swallowing or lying flat.

\section*{Risks and Limitations}
Transthoracic echocardiography is noninvasive. Transesophageal echocardiography carries sedation and throat or esophageal-injury risks. Image quality and the type of study limit what can be assessed.

\section*{Results and Interpretation}
Results may describe chamber size, pumping function, valve structure, blood-flow gradients, and pressure estimates. The clinician interprets these findings with symptoms, examination findings, and other tests.

\section*{Aftercare and Recovery}
After a standard transthoracic study, usual activities generally resume immediately. After sedation or a transesophageal study, follow the facility's restrictions and arrange transportation as instructed.

\section*{When to Seek Medical Care}
Follow the procedure team's specific instructions. Seek urgent medical assessment for severe or worsening pain, uncontrolled bleeding, fever or signs of infection, trouble breathing, chest pain, fainting, new weakness or confusion, inability to urinate, or any rapidly worsening symptom. The appropriate warning signs depend on the procedure and the patient's medical condition.

\section*{References}
\begin{enumerate}
  \item MedlinePlus. Echocardiogram. U.S. National Library of Medicine; 2025.
  \item Current Medical Diagnosis and Treatment. McGraw-Hill; 2025.
\end{enumerate}
\clearpage
